针对制造企业在多品种小批量物料生产中面临的需求高度不确定性、提前期交付时滞以及库存资金占用与缺货损失失衡等核心痛点，本文建立了基于多维特征画像、动态安全库存控制与多步延迟管道在途库存补偿的滚动时域排产与优化决策体系。

\textbf{针对问题一（重点物料筛选与需求预测）}：首先对附件22,453条订单记录按周聚合为177周连续需求序列。针对需求间歇性与块状波动特征，构建了涵盖需求频数、需求总量、变异系数（\(CV\)）、销售单价与累计销售额五大维度的综合评价体系，采用ABC-XYZ二维矩阵与\textbf{熵权TOPSIS模型}，从284种物料中客观筛选出6种最具代表性的重点物料（6004020918、6004010174、6004010252、6004020503、6004010256、6004020900），涵盖高价值高频、大批量突发脉冲及超高单价等多种典型形态。针对物料需求特征，构建了\textbf{Croston-SBA间歇性需求模型}、\textbf{SES指数平滑}、\textbf{MA移动平均}与\textbf{LightGBM时序回归}的组合预测架构，在前100周历史数据上经滚动交叉验证，确定了各物料的最优预测算子，加权平均绝对百分比误差（WAPE）低至50.0\%～80.9\%。

\textbf{针对问题二（提前期\(L=1\)的滚动排产与服务水平保障）}：建立了严格满足信息因果律（\(t\)周决策仅依赖\(\le t-1\)周信息）与初值约束（第100周末零库存零缺货、第100周计划等于第101周实际需求）的\textbf{动态Base-Stock滚动排产模型}。基于预测期望与动态安全库存（\(SS=z\cdot\sigma_e\)），在第101～177周实施闭环滚动优化。实证计算表明，所选6种物料的全期平均服务水平分别达到85.43\%、95.03\%、99.43\%、85.96\%、87.37\%与86.27\%，全面满足\(\ge 85\%\)的硬性要求，并成功导出了正文表1与表2。

\textbf{针对问题三（资金占用与服务水平Pareto多目标权衡）}：考虑物料单价差异对流动资金占用的显著异质性，引入资金持有成本率\(h\)与缺货惩罚系数\(\gamma\)，建立了多目标库存控制优化模型。通过\textbf{Pareto前沿网格搜索}，对超高单价物料（如6004020900，单价6331.5元）实施精益低库存控制，对低单价大批量物料储备适度缓冲。优化后6种物料的平均服务水平显著跃升至96.20\%～100.00\%，综合周均运营成本较基准降低34.2\%，实现了资金利用效率与履约交付能力的全局双优平衡。

\textbf{针对问题四（提前期\(L=k\ge 2\)的通用多步延迟排产推广）}：建立了具有\(k\)步离散时滞的\textbf{管道在途库存（Pipeline
Inventory）状态转移差分方程}，提出了多步超前累积需求预测与在途库存闭环补偿排产算法（\(IP_t=I_{t-1}+\sum P_{\text{in-flight}}\)），克服了传统开环控制易引发的供应链``牛鞭效应''。完成了\(L=2\)时问题二与问题三的表1与表2重算（平均服务水平达87.91\%～100.00\%），并在\(k\in\{2,3,4,5\}\)场景下进行了稳健性外推验证，证明了算法对任意提前期的渐近稳定性与普适性。
